\label{sec:results}

% Figures are written to <root>/figures by src/main.py and src/evaluate.py.
% Copy the ones cited below into report/figures/ before compiling, since
% \includegraphics resolves relative to this directory.
% Numbers marked in red are filled from results/metrics.csv once the runs finish.

% \pending draws the figure once a run has produced it, and a labelled box
% until then, so the report compiles at every stage of the experiments.
\providecommand{\pending}[2][\textwidth]{%
    \IfFileExists{#2}{\includegraphics[width=#1]{#2}}{%
        \fbox{\parbox[c][3cm][c]{\dimexpr#1-2\fboxsep-2\fboxrule\relax}{%
            \centering\footnotesize\texttt{\detokenize{#2}}\\[4pt]
            \normalfont\textcolor{red}{awaiting the run}}}%
    }%
}

% Due to the challenges inherent in the dataset, we do not expect \texttt{X3D} to perform comparably to purpose–built sign language models, but rather to reach an acceptable accuracy while retaining an inference cost small enough for edge deployment.
% Following \textbf{Sec. \ref{sec:introduction}}, we are further interested in how far fine–tuning depth affects that trade–off.

\subsubsection{Learning behaviour on train/validation sets}

We train every combination of the three variants \texttt{X3D\_XS}, \texttt{X3D\_S} and \texttt{X3D\_M} with $0$ to $5$ of the six blocks held frozen.
\textbf{Figure \ref{fig:curves}} shows the two extremes of this sweep, and \textbf{Table \ref{tab:val_sweep}} collects the best validation top--1 of every run.

\begin{figure}[H]
    \centering
    \begin{subfigure}[b]{0.49\textwidth}
        \pending{figures/learning/x3d_xs_freeze5.png}
        \caption{\texttt{X3D\_XS}, head only.}
    \end{subfigure}
    \hfill
    \begin{subfigure}[b]{0.49\textwidth}
        \pending{figures/learning/x3d_xs_freeze0.png}
        \caption{\texttt{X3D\_XS}, full fine--tune.}
    \end{subfigure}

    \vspace{0.4em}

    \begin{subfigure}[b]{0.49\textwidth}
        \pending{figures/learning/x3d_s_freeze5.png}
        \caption{\texttt{X3D\_S}, head only.}
    \end{subfigure}
    \hfill
    \begin{subfigure}[b]{0.49\textwidth}
        \pending{figures/learning/x3d_s_freeze0.png}
        \caption{\texttt{X3D\_S}, full fine--tune.}
    \end{subfigure}

    \vspace{0.4em}

    \begin{subfigure}[b]{0.49\textwidth}
        \pending{figures/learning/x3d_m_freeze5.png}
        \caption{\texttt{X3D\_M}, head only.}
    \end{subfigure}
    \hfill
    \begin{subfigure}[b]{0.49\textwidth}
        \pending{figures/learning/x3d_m_freeze0.png}
        \caption{\texttt{X3D\_M}, full fine--tune.}
    \end{subfigure}
    \caption{Per--epoch cross--entropy loss and top--1 accuracy on the train and validation splits, for the two extreme fine--tuning depths of each variant. Freezing five blocks leaves only the classifier trainable, whereas freezing none updates the whole network.}
    \label{fig:curves}
\end{figure}

\begin{table}[H]
\centering\small
\begin{tabular}{@{}c cc cc cc@{}}
\toprule
\textbf{Frozen} & \multicolumn{2}{c}{\texttt{X3D\_XS}} & \multicolumn{2}{c}{\texttt{X3D\_S}} & \multicolumn{2}{c}{\texttt{X3D\_M}} \\
\cmidrule(lr){2-3}\cmidrule(lr){4-5}\cmidrule(lr){6-7}
\textbf{blocks} & train & val & train & val & train & val \\
\midrule
$0$ & $95.4$ & $63.9$ & $98.3$ & $\mathbf{81.5}$ & $98.4$ & $85.4$ \\
$1$ & $95.4$ & $\mathbf{64.7}$ & $98.4$ & $80.4$ & $98.3$ & $85.5$ \\
$2$ & $95.5$ & $64.2$ & $98.4$ & $81.0$ & $98.9$ & $\mathbf{86.1}$ \\
$3$ & $93.7$ & $58.7$ & $98.5$ & $75.8$ & $98.7$ & $83.4$ \\
$4$ & $92.4$ & $49.2$ & $97.7$ & $71.1$ & $97.8$ & $77.1$ \\
$5$ & $48.0$ & $19.9$ & $64.3$ & $36.5$ & $69.8$ & $38.5$ \\
\bottomrule
\end{tabular}
\caption{Top--1 accuracy (\%) on the train and validation splits, for each variant and freezing depth. The variants differ only in the clip they sample, i.e. \texttt{X3D\_XS} takes $4$ frames at $160^2$, \texttt{X3D\_S} $13$ frames at $160^2$ and \texttt{X3D\_M} $16$ frames at $224^2$. The best validation entry per variant is bolded.}
\label{tab:val_sweep}
\end{table}

\paragraph{Underfitting.} Training the classifier head only amount to underfitting, evidently $48.0\%$, $64.3\%$ and $69.8\%$ accuracy for \texttt{X3D\_XS}, \texttt{X3D\_S} and \texttt{X3D\_M}, respectively.
This is indicative of the \texttt{Kinetics–400} pretrained weights may not contribute vastly to models' performance on \texttt{AUTSL}, and mostly rely on the learning capacity of the classifier block.

\paragraph{Pretrained weights may not help.} Looking at learning curves in \textbf{Figure \ref{fig:curves}}, at the first few epochs, they accuracy is extremely low.
Semantically, this is true, since \texttt{Kinetics–400} and \texttt{AUTSL} are constructed for two completely different tasks. Nonetheless, we cannot strongly claim that this does not help, since the pretrained weights may yield a relatively good initial starting point to optimise the models, in terms of local optimality and convergence speed.
Further ablation studies must be conducted to compare between train–from–scratch vs. train–from–pretrained.

\paragraph{Overfitting.} Interestingly, enabling more blocks in the backbone model trainable flips the situation. It immediately overfits to the training data for the first block adjacent to the classifier head, evidently by the huge gap between train and validation accuracy.
Nevertheless, the gap is converging as we allows more parameters learnable, and astonishingly, such train–validation gap is narrower in \texttt{X3D\_M} than in \texttt{X3D\_XS}.

\paragraph{Clip sampling outweighs adaptation depth.} The ordering $\texttt{X3D\_M} > \texttt{X3D\_S} > \texttt{X3D\_XS}$ holds at every freezing depth, and the margin is large.
More specifically, on validation accuracy, \texttt{X3D\_M} with three blocks frozen ($83.4\%$), fully fine–tuned ``best'' \texttt{X3D\_S} ($81.5\%$), and one block frozen ``best'' \texttt{X3D\_XS} ($64.7\%$).
From \textbf{Table \ref{tab:main_results}}, three variants differ only in the clip drawn from each video and the cropping size each frame.
Therefore, the gain is attribute to temporal and spatial sampling rather than to capacity, and it is traded off with GFLOPs.

\paragraph{Recommendation.} From the above–mentioned results and analysis, we shall pick \texttt{X3D\_M} with 2 frozen blocks as our ``best'' model, due to its performance competency and good trade–off with trainable parameters.

\subsubsection{Final evaluation on test set}

\begin{table}[H]
\centering\small
\begin{tabular}{@{}l c c c c c c c@{}}
\toprule
\textbf{Model} & \textbf{Frozen} & \textbf{Top--1} & \textbf{Top--3} & \textbf{Top--5} & \textbf{Params} & \textbf{Trainable} & \textbf{GFLOPs} \\
 & & (\%) & (\%) & (\%) & (M) & (M) & \\
\midrule
\texttt{X3D\_XS} & $0$ & $54.08$ & $72.36$ & $78.40$ & $3.438$ & $3.438$ & $0.63$ \\
\texttt{X3D\_XS} & $1$ & $\mathbf{56.91}$ & $74.58$ & $80.78$ & $3.438$ & $3.437$ & $0.63$ \\
\texttt{X3D\_XS} & $2$ & $53.49$ & $73.08$ & $79.76$ & $3.438$ & $3.422$ & $0.63$ \\
\texttt{X3D\_XS} & $3$ & $49.67$ & $69.31$ & $76.90$ & $3.438$ & $3.348$ & $0.63$ \\
\texttt{X3D\_XS} & $4$ & $38.89$ & $58.49$ & $66.61$ & $3.438$ & $2.779$ & $0.63$ \\
\texttt{X3D\_XS} & $5$ & $14.68$ & $26.01$ & $31.36$ & $3.438$ & $1.432$ & $0.63$ \\
\midrule
\texttt{X3D\_S} & $0$ & $75.67$ & $89.74$ & $93.42$ & $3.438$ & $3.438$ & $2.03$ \\
\texttt{X3D\_S} & $1$ & $\mathbf{76.56}$ & $89.76$ & $93.26$ & $3.438$ & $3.437$ & $2.03$ \\
\texttt{X3D\_S} & $2$ & $75.67$ & $89.60$ & $93.50$ & $3.438$ & $3.422$ & $2.03$ \\
\texttt{X3D\_S} & $3$ & $69.85$ & $85.89$ & $90.35$ & $3.438$ & $3.348$ & $2.03$ \\
\texttt{X3D\_S} & $4$ & $63.46$ & $81.05$ & $86.93$ & $3.438$ & $2.779$ & $2.03$ \\
\texttt{X3D\_S} & $5$ & $29.35$ & $43.89$ & $50.12$ & $3.438$ & $1.432$ & $2.03$ \\
\midrule
\texttt{X3D\_M} & $0$ & $\mathbf{82.38}$ & $93.50$ & $95.94$ & $3.438$ & $3.438$ & $4.90$ \\
\texttt{X3D\_M} & $1$ & $80.83$ & $92.70$ & $95.64$ & $3.438$ & $3.437$ & $4.90$ \\
\texttt{X3D\_M} & $2$ & $82.28$ & $94.49$ & $96.52$ & $3.438$ & $3.422$ & $4.90$ \\
\texttt{X3D\_M} & $3$ & $77.17$ & $90.56$ & $94.04$ & $3.438$ & $3.348$ & $4.90$ \\
\texttt{X3D\_M} & $4$ & $73.62$ & $89.20$ & $93.40$ & $3.438$ & $2.779$ & $4.90$ \\
\texttt{X3D\_M} & $5$ & $34.86$ & $50.98$ & $57.31$ & $3.438$ & $1.432$ & $4.90$ \\
\bottomrule
\end{tabular}
\caption{Performance on the \texttt{AUTSL} test split ($3{,}741$ clips evaluated over $226$ classes). The best top–1 per variant is bolded.}
\label{tab:main_results}
\end{table}

\paragraph{Freezing depth.} Training the whole network rather than the head alone adds $39.4$, $46.3$ and $47.5$ top--1 points for \texttt{X3D\_XS}, \texttt{X3D\_S} and \texttt{X3D\_M}, at $2.4\times$ the trainable parameters.
Interestingly, by unfreezing only one block adjacent to the head block, i.e. \texttt{res}$_5$, lifts \texttt{X3D\_M} from $34.86\%$ to $73.62\%$ for $1.35$M extra trainable parameters, likewise for \texttt{X3D\_XS} and \texttt{X3D\_S}.
Thereafter, unfreezing other blocks gain more accuracy, though not a big jump as before.

\paragraph{Model size.} Moving from \texttt{X3D\_XS} to \texttt{X3D\_S} costs $3.2\times$ the GFLOPs and returns $19.7$ top--1 points; the further step to \texttt{X3D\_M} costs $2.4\times$ again for only $5.8$.
The longer clip is therefore not wasted on isolated signs, but its value saturates.

\paragraph{Error type.} The top--1 to top--3 margin is widest for \texttt{X3D\_XS} and narrowest for \texttt{X3D\_M}.
We shall claim, though insignificantly, that the small variant's errors are largely near misses among confusable signs, whereas the residual errors of \texttt{X3D\_M} are more often outright failures.

\subsubsection{Generalising to our data}

We separate this section from the testing above, since the test split still shares underlying patterns with the train split, e.g. signer identity, background and experimental bias.
To build the dataset we watched the public clips, replicated their motions and annotated with the same numeric \texttt{id}, as the model learnt indexed rather than text--based classes.
Also, it poses more challenges, as we cannot mimic the mouth movements that some signs involve, our lighting is poor and introduces scattering, and \texttt{AUTSL} records the full body whereas we recorded only the upper half.

\paragraph{Setup.} We recorded ten clips on a laptop camera at $1620\times1080$, each imitating one \texttt{AUTSL} sign and named after the test clip whose label it borrows.
In other words, the footage itself is our own and only the labelling convention is shared.
They were converted to \texttt{mp4} and passed through the pipeline of \textbf{Sec. \ref{sec:implementation}}. We run the inference on this hand–crafted data using CPU on MacBook M3 Pro.

\begin{table}[H]
\centering\small
\begin{tabular}{@{}c c c c c@{}}
\toprule
\textbf{Clip} & \textbf{True class} & \textbf{Predicted (top–1)} & \textbf{Confidence} & \textbf{Inference time (s)} \\
\midrule
$0$ & $161$ & $161$ & $1.000$ & 0.369 \\
$1$ & $77$  & $77$  & $0.667$ & 0.331\\
$2$ & $22$  & $18$  & $0.974$ & 0.343\\
$3$ & $157$ & $213$ & $0.860$ & 0.310\\
$4$ & $112$ & $32$  & $0.680$ & 0.328\\
$5$ & $128$ & $54$  & $0.703$ & 0.313\\
$6$ & $14$  & $14$  & $1.000$ & 0.316\\
$7$ & $5$   & $5$   & $1.000$ & 0.340\\
$8$ & $23$  & $23$  & $0.987$ & 0.332\\
$9$ & $155$ & $155$ & $0.999$ & 0.404\\
\bottomrule
\end{tabular}
\caption{Predictions of \texttt{X3D\_M} with two frozen blocks on ten self--recorded clips. Confidence is the \texttt{softmax} probability of the predicted class.}
\label{tab:custom}
\end{table}

\paragraph{Accuracy.} As collected in \textbf{Table \ref{tab:custom}}, the model recognises six of the ten clips, i.e. $60\%$ top--1, $60\%$ top--3 and $70\%$ top--5, yet incomparable with the test results in \textbf{Table \ref{tab:main_results}} due to insufficient samples.

\paragraph{Confidence.} The wrong predictions are as confident as the right ones, i.e. $0.68$ to $0.97$ against $0.67$ to $1.00$.
Softmax confidence therefore offers no usable signal for rejecting an out--of--distribution clip.

\paragraph{Latency.} Inference took $0.310$ to $0.404$ seconds per clip, $0.339$ on average, on a CPU alone.
It suggests the model is deployable for real–time inference.
